% Fresh reconstruction after the 09:11 UTC environment replacement.
\subsection{Quadratic representation growth does not force generation}

\begin{example}[Proved: an alternating-product obstruction in every coefficient characteristic]
\label{D:quadratic-growth-alternating-obstruction}
For odd integers \(n\geq5\), put \(m_n=2^{n^2}\), and let
\[
 \Pi=\prod_{\substack{n\geq5\\n\ {\rm odd}}} A_n^{\,m_n}.
\]
This profinite group has the following properties.
\begin{enumerate}
\item It is small: for each finite group \(H\), the set of continuous
homomorphisms \(\Pi\to H\) is finite.
\item Every open subgroup has finite maximal prosolvable quotient.
\item For every finite field \(\mathbf F_q\), of any characteristic,
and every finite-dimensional continuous \(\mathbf F_q[\Pi]\)-module \(V\),
\[
 \dim_{\mathbf F_q}H^1(\Pi,V)\leq2\dim_{\mathbf F_q}V.
 \tag{A1}
\]
\item There is one constant \(C\), independent of the characteristic,
\(q\), and \(r\), such that
\[
 \#\bigl(
 \operatorname{Hom}_{\mathrm{cts}}(\Pi,\operatorname{GL}_r(\mathbf F_q))
             /\operatorname{GL}_r(\mathbf F_q)\bigr)
       \leq q^{Cr^2}.
 \tag{A2}
\]
\end{enumerate}
Nevertheless \(\Pi\) is not topologically finitely generated.
This is an abstract profinite obstruction. No realization as a
geometric fundamental group is asserted.
\end{example}

\begin{proof}
We give the counting argument in all coefficient characteristics.
We use the simplicity of \(A_n\) for \(n\geq5\), with its elementary
proof in \cite[Theorem 4.33 and Lemmas 4.35--4.36, pp.~69--70]
{MilneGT2025B}. For odd \(n\), \(A_n\) is generated by the even
\(n\)-cycle \((1\,2\,\cdots\,n)\) and \((1\,2\,3)\).
Indeed their conjugates include the consecutive three-cycles.
Those generate \(A_n\): inductively they contain \(A_{n-1}\), and
a three-cycle involving the last letter, together with its
\(A_{n-1}\)-conjugates, supplies every three-cycle involving that
letter. The initial cases are immediate. In particular every \(A_n\)
is perfect and two-generated.

\par\smallskip\noindent\textbf{A dimension bound for nontrivial factors.}\quad
Let \(W\) be a nontrivial representation of \(A_n\) over an
algebraically closed field, and put \(a=\dim W\).
Simplicity makes the representation faithful.
In characteristic different from two, use the subgroup
\((C_2)^{\lfloor n/4\rfloor}\) generated by disjoint double
transpositions. In characteristic two, use
\((C_3)^{\lfloor n/3\rfloor}\) generated by disjoint three-cycles.
These groups have order prime to the coefficient characteristic,
so their actions are simultaneously diagonalizable.
For an elementary abelian group of rank \(s\), a faithful diagonal
representation needs at least \(s\) characters: their common
kernel is zero only if they span its character space.
It follows in every characteristic that
\[
 n\leq8a,\qquad a\geq2.
 \tag{A3}
\]
The last inequality also follows from perfectness.

\par\smallskip\noindent\textbf{Absolutely simple modules.}\quad
A continuous finite-dimensional module factors through a finite
product of the coordinate groups. Over an algebraically closed
field, the simple modules for a product are the outer tensor
products of the simple modules of its factors, including in
modular characteristic. To see this without a semisimplicity
assumption on group algebras, take their finite-dimensional
algebras. The ideal generated by their radicals is nilpotent in
the tensor product, while the quotient is a product of matrix
algebras. Its simple modules are exactly the stated tensor
products. We use only these standard semisimple quotient facts
\cite[Theorems 7.25, 7.29, 7.36--7.37, pp.~110--114]{MilneGT2025B}.

Let a simple module have nontrivial dimensions
\(a_1,\ldots,a_t\) at its support coordinates, whose alternating
degrees are \(n_1,\ldots,n_t\). Its dimension is
\(r=\prod_i a_i\). For \(a_i\geq2\),
\[
 t\leq\log_2r,\qquad
 \sum_i a_i^2\leq\prod_i a_i^2=r^2,\qquad
 \sum_i n_i^2\leq64r^2.
 \tag{A4}
\]
The middle inequality follows by induction from
\(a^2+b^2\leq a^2b^2\) for \(a,b\geq2\).
There are at most \((8r)^{1+\lfloor\log_2r\rfloor}\) ordered choices
of the \(n_i\); this overcounts support types and is sufficient.
For each such support, its coordinate-copy choices number at most
\[
 \prod_i m_{n_i}=2^{\sum_i n_i^2}.
\]
For \(A_n\), the number of simple module labels is at most
\(\dim k[A_n]=|A_n|\leq n^n\leq2^{n^2}\).
Thus the label choices have the same bound.
Combining these estimates shows that the number of continuous
absolutely simple \(r\)-dimensional modules is at most
\[
 2^{C_0r^2}
 \tag{A5}
\]
for an absolute constant \(C_0\), uniformly in characteristic.
The rank-one trivial module is included by increasing the constant.

Over \(\mathbf F_q\), simple modules are classified by Frobenius
orbits of the absolutely simple modules, with unique descent.
Here is an algebraic justification valid also in modular
characteristic. A finite-dimensional semisimple algebra over
a finite field is a product of matrix algebras over finite
extensions of that field, since finite division rings are
commutative. A simple module for
\(\operatorname{Mat}_a(\mathbf F_{q^e})\) becomes, after algebraic
closure, the direct sum of the \(e\) distinct Frobenius-conjugate
absolutely simple modules of dimension \(a\). This description
also proves uniqueness. Apply it to the semisimple quotient of
the finite group algebra through which the module factors.
Consequently the number \(s_r(q)\) of simple modules of dimension
\(r\) over \(\mathbf F_q\) satisfies
\[
 s_r(q)\leq\sum_{e\mid r}2^{C_0(r/e)^2}
             \leq q^{C_1r^2},
 \tag{A6}
\]
with \(C_1\) absolute and all \(q\geq2\).

\par\smallskip\noindent\textbf{Uniform extension bounds.}\quad
For two absolutely simple modules \(U,V\), the degree-one
Kunneth formula for a finite product gives
\[
 \operatorname{Ext}^1(U,V)=
 \bigoplus_i
 H^1\!\left(A_{n_i},\operatorname{Hom}(U_i,V_i)\right)
       \otimes\!\!\bigotimes_{j\ne i}
                  \operatorname{Hom}_{A_{n_j}}(U_j,V_j).
 \tag{A7}
\]
One obtains this formula by tensoring the ordinary free bar
resolutions of the factors. Over a field, the degree-zero and
degree-one cohomology of the resulting total complex have exactly
the displayed decomposition. There is no restriction on the
coefficient characteristic.

For a two-generated finite group and a module \(W\), a derivation
is determined by its two values on the generators. Hence
\(\dim H^1(G,W)\leq2\dim W\).
By Schur's lemma, the invariant Hom spaces in (A7) are zero
unless the two simple factors agree, and are one-dimensional
when they do agree. If \(U,V\) differ at two support factors,
(A7) is zero. If they differ at exactly one, it is bounded by
\(2\dim U\dim V\). If \(U=V\), its dimension is at most
\(2\sum_i a_i^2\leq2(\dim U)^2\), by (A4).
Coordinates trivial on both sides contribute
\(H^1(A_n,k)=0\), because \(A_n\) is perfect.

These computations do calculate extensions for the full infinite
product. In an extension of two modules trivial on a coordinate
factor, that factor acts through the additive group
\(\operatorname{Hom}(U,V)\); perfectness kills this action.
Thus no new support coordinate can occur in the extension.
All groups in (A7) may be computed on a common finite product.

Cohomology commutes with extension of a coefficient field on this
finite product. Splitting the finite-field simple modules over
an algebraic closure and summing the preceding bounds therefore
gives
\[
 \dim_{\mathbf F_q}\operatorname{Ext}^1(U,V)
       \leq2\dim U\dim V
 \tag{A8}
\]
for simple finite-field modules. Devissage in both variables,
using the long exact Ext sequences, gives the same inequality
for arbitrary finite-dimensional modules. Taking \(U\) trivial
proves (A1).

To count all \(r\)-dimensional modules, fix an ordered composition
\(r=r_1+\cdots+r_t\). There are at most
\(q^{C_1\sum_i r_i^2}\) choices of the simple factors.
Successively extending them gives at most
\[
 q^{\,2\sum_{i<j}r_ir_j}
\]
choices, by (A8). We count extension classes with fixed
identifications of their ends, which can only overcount
isomorphism classes of middle terms. Since
\[
 C_1\sum_i r_i^2+2\sum_{i<j}r_ir_j
 \leq \max(C_1,1)r^2,
\]
and there are \(2^{r-1}\) ordered compositions of \(r\),
(A2) follows with one larger absolute constant.

\par\smallskip\noindent\textbf{Smallness and failure of generation.}\quad
A homomorphism from \(\Pi\) to a finite group \(H\) is trivial on
every \(A_n\) with \(|A_n|>|H|\), by simplicity. Only finitely many
coordinate factors remain, so there are finitely many such
homomorphisms. This proves smallness.

An open subgroup contains the full tail product after finitely
many coordinates. That tail has no nontrivial finite solvable
quotient, since each of its factors is perfect. Every solvable
quotient of the open subgroup therefore factors through a fixed
finite head quotient. Its maximal prosolvable quotient is finite.

Finally suppose \(s\) elements generated \(\Pi\).
Their projection to \(A_n^{m_n}\) would generate that finite group.
The \(s\)-tuples at any two coordinates must be different: if
they were equal, the generated subgroup would have equal entries
in those two coordinates and could not be the full product.
Thus \(m_n\leq|A_n|^s\).
But \(2^{n^2}>(n!)^s\) for all sufficiently large odd \(n\), since
\(n\log 2-s\log n\to+\infty\).
This contradiction proves that \(\Pi\) is not finitely generated.
\end{proof}

\begin{remark}[What the obstruction excludes]
Even the conjunction of smallness, finite maximal prosolvable quotients
for every open subgroup, all-characteristic rank-linear \(H^1\), and
a uniform all-rank quadratic representation-class bound does not imply
finite generation. A geometric proof must impose a further condition,
for example a sufficiently strong bound on nonabelian simple
multiplicities. The present example is not a counterexample to any
of the geometric statements under investigation.
\end{remark}

